\section{Free groups, equation languages and tree actions}\label{sec:free}

\subsection{Growth of a nonprimitive automorphic orbit (F41)}\label{sec:f41}
\entry{F41}{D.~Puder and C.~Wu}
For a nonprimitive word in $F_r$, is the automorphic orbit's exponential
growth rate $\sqrt{2r-1}$? The identity is an exception to the literal
wording: its orbit has ball growth one. We use ordinary length in a
fixed basis; this is not a stronger word-dependent cyclic-orbit conjecture.

\begin{theorem}\label{prop:f41}
Let $r\ge2$, $q=2r-1$, and let $w\ne1$ be nonprimitive in $F_r$.
There are constants $C_w>0$, $D_w\ge0$ such that
\[
 |\Aut(F_r)w\cap B(n)|\le C_w(n+1)^{D_w}q^{n/2}.
\]
The ball growth limit and the spherical growth limsup are $\sqrt q$.
\end{theorem}

Two imported results are essential. Kharlampovich--Myasnikov's
arbitrary-definable-set theorem says that a definable subset of $F_r$
or its complement is a sub-multipattern~\cite[Definitions 5--6,
Theorem 13]{km}. A one-variable pattern is described by finitely many
noncancelling equations with fixed coefficients; every variable in its
first equation is a repeated piece of the output, possibly with inverse
orientation and overlap. A sub-multipattern is a subset of a finite
union of these patterns. Pillay's parameter-free generic type $p_0$
is realized by basis elements, and its realizations in finite-rank free
groups are primitive~\cite[Fact 1.10, Theorem 2.1]{pillay}. Every
parameter-free definable set containing a generic element has finitely
many left translates covering the whole group.

\begin{proof}[Derivation of Proposition~\ref{prop:f41}]
First count reduced words with a fixed cover by repeated pieces. Suppose
a first pattern equation has $t$ variable blocks and $C$ coefficient
letters. For length $n$, choose the block lengths and a second occurrence
of each block, including orientation. There are at most a constant
times $(n+1)^{2t}$ choices. In each choice form the signed graph of
identified letter positions. An inconsistent inverse cycle makes the
choice impossible. In a consistent choice every noncoefficient position
is joined to a different position, so the number $c$ of components
satisfies $c\le(n+C)/2$.

Collapse these components in the path of consecutive positions. The
quotient is connected. On a spanning tree, the root component has
$2r$ possible letters, while each further component has at most $q$:
one letter is excluded by free reduction across its parent edge, with
the prescribed signs. Ignoring other constraints only enlarges the
count. Thus each choice contributes at most $2r q^{(n+C)/2}$ words.
Finite unions and subsets preserve a polynomial times $q^{n/2}$ bound.
No sub-multipattern is finite-translate generic, since translating by
a fixed element changes the radius by at most a constant, whereas
$F_r$ has ball growth $q$.

Since $w$ is nonprimitive, it does not realize $p_0$. Choose a
parameter-free formula $\varphi\in p_0$ false at $w$. Its negation
defines an automorphism-invariant set $P$ containing the orbit of $w$.
The complement contains a basis element and is finite-translate generic.
The definable-set dichotomy and the preceding count therefore make $P$,
not its complement, a sub-multipattern. This proves the upper bound.
No definability of the orbit or of the primitive set is assumed.

For the lower bound, replace $w$ by a cyclically reduced conjugate $v$,
of length $\ell$. There are at least $(2r-2)q^{m-1}$ reduced words
$t$ of length $m$ for which $tvt^{-1}$ is reduced as written: exclude
the inverse of the first letter of $v$ and the last letter of $v$ from
the last letter of $t$. Distinct $t$ give distinct words of length
$2m+\ell$. This gives the ball lower bound and the matching spherical
limsup. Empty parity classes can prevent a spherical limit.
\end{proof}

In rank one every orbit is finite. The higher-rank proof is a proposed
consequence of the two deep cited theorems and the piece-cover count;
its finite checks do not verify the model-theoretic interface.
\evidence{F41}{multipattern-proof.md}

\subsection{Discriminating a free square (GA3)}\label{sec:ga3}
\entry{GA3}{A.~Miasnikov}
If $G$ acts freely on a $\Lambda$-tree, can each finite set of
nonidentity elements of $G*G$ survive a homomorphism to $G$?
For nontrivial abelian $G$ the literal answer is negative: a commutator
between the two factors always dies. Free actions below have no inversions.

\begin{theorem}\label{prop:ga3}
If $G$ is nonabelian and $\Lambda$-free, for each finite
$W\subset(G*G)\setminus\{1\}$ there are $h\in G$ and $N$ such
that the maps which are the identity on the first factor and
$b\mapsto h^nbh^{-n}$ on the second are nontrivial on $W$ for all
$n\ge N$. No finite-generation hypothesis on $G$ is required.
\end{theorem}
\begin{proof}[Derivation]
It suffices first to control a finitely generated nonabelian subgroup
$K$. In a divisible scalar extension of its free $\Lambda$-tree,
choose a finite symmetric generating set $S$ and the scale
\[
 L=\max\{\ell(s),\ell(st):s,t\in S\}>0.
\]
The disjoint-axis formula bounds pairwise distances between generator
axes by $L/2$. The finite Helly property gives a point within $L$
of every axis, hence with generator displacements at most $3L$.
The hull of its orbit has distances in the principal convex subgroup
generated by $L$. Quotient by its maximal infinitesimal convex subgroup
and extend the resulting Archimedean ordered group to $\R$.
An element attaining translation length $L$ stays hyperbolic, so this
collapse does not have a global fixed point.

Guirardel's collapse theorem and its general-$\Lambda$ remark give
abelian arc fixators and trivial tripod fixators~\cite[Fact 5.1]{guirardel}.
The final scalar extension preserves these properties because the
divisible Archimedean subgroup is dense and no branch points are added.
The CSA property~\cite[Lemma 2.5]{guirardel} implies that a nonabelian
$K$ has no nontrivial abelian normal subgroup: such a subgroup would
lie in every conjugate of its maximal abelian subgroup, contradicting
malnormality. Thus $K$ preserves no line and fixes no end. For a line,
the arc kernel is abelian and the remaining faithful torsion-free action
is by translations. For an end, the Busemann kernel is abelian because
any two of its elements fix a common terminal ray. Either case would
make $K$ abelian.

In this real-tree action each nonidentity element fixes at most two
ends, since three fixed ends give a fixed tripod. Given finitely many
nonidentity coefficients $C\subset K$, choose two limit-set ends
whose pair is disjoint from all its $C$-translates. Tree ping-pong,
or the elementary separator in the retained supplement, gives a
hyperbolic $h$ with both ends in small prescribed neighborhoods
$U_+,U_-$ and
\[
 c(U_+\cup U_-)\cap(U_+\cup U_-)=\varnothing\quad(c\in C).
\]
Sufficient powers send the complement of $U_-$ into $U_+$ and the
complement of $U_+$ into $U_-$ in the reverse direction. This is the
standard boundary separation mechanism; its prior use is credited
in the \record{problems/GA3/audit.md}{source audit}.

Cyclically reduce each word in $W$. A core in one free factor survives
automatically. Write every other core as $a_1b_1\cdots a_mb_m$,
and include all coefficients, together with two noncommuting elements
of $G$, in the generators of $K$. Starting with an end in $U_-$,
read the substituted core from right to left. Each $h^{-n}$, $b_i$,
$h^n$, $a_i$ sends it successively to $U_-$, outside both neighborhoods,
$U_+$ and outside both. Its final image is outside $U_-$, so the
word is nonidentity. The same $h,N$ work for all cores. The factor maps
are automorphisms of all of $G$, which extends the argument from
$K*K$ to $G*G$.
\end{proof}

The trivial group satisfies the property vacuously. This separates the
literal abelian obstruction from the substantive nonabelian assertion.
The arbitrary-$\Lambda$ collapse, not CSA alone, is the principal
structural dependency. No procedure for obtaining $h,N$ from an
unspecified presentation is asserted.
\evidence{GA3}{proof.md}

\subsection{A shared polynomial-identity method}\label{sec:languages}
Let $F_r$ have basis $A$, and let $B_\phi$ be the exponent-sum matrix
of an endomorphism. If $\det B_\phi\ne0$, then $\phi$ is injective:
its image is a free group generated by $r$ elements, and its ambient
abelianization has rank $r$, so its rank is exactly $r$; Hopficity then
applies to the surjection onto that image.

We use the effective \emph{full solution relation} of word equations
with involution and regular constraints~\cite[Theorem 4.3]{diekert-elder}.
It has the form
\begin{equation}\label{eq:edt0l}
 \{(h(d_1),\ldots,h(d_k)):h\in\mathcal R\},
\end{equation}
where $\mathcal R$ is a regular control of explicit endomorphisms of
one finite extended alphabet, with separately designated output letters
$d_i$. Group equations reduce to these monoid equations by triangulation.
For reduced $X,Y,Z$, a triangle $XY=Z$ is replaced by
\[
 X=ST,\qquad Y=\overline T R,\qquad Z=SR,
\]
with reduced-word constraints; maximal cancellation proves equality of
the projected solution relations. Cyclic reduction and positivity are
regular constraints in the prescribed basis. Recognizing monoids can
be made compatible with involution using $M\times M^{\mathrm{op}}$.

\begin{lemma}\label{lem:polynomial}
One can decide whether a rational polynomial in the separate output
letter counts vanishes on every tuple of~\eqref{eq:edt0l}.
\end{lemma}
\begin{proof}
Incidence matrices act linearly on the stacked count vector $y\in\N^D$.
For polynomial degree $d$, lift $y$ to the vector $\nu_d(y)$ of all
monomials of degree at most $d$, including one. Its dimension is
$\binom{D+d}{d}$. Each transition induces an effective rational matrix
$T^{[d]}$ satisfying $\nu_d(Ty)=T^{[d]}\nu_d(y)$.
In each automaton state compute the rational span of reachable lifted
vectors. Propagating these spaces can strictly increase total dimension
only finitely many times. The stable spaces are exactly the spans of
actual path vectors, by induction in both directions. Test the linear
functional representing the polynomial on all accepting spaces.
Reverse the control automaton if required by morphism-composition order.
If terminal filtering is not built into the chosen convention, include
the finite support of the count vector in the control state. This also
handles erasing morphisms.
\end{proof}

This decides a universal identity, equivalently the existence of a
\emph{nonzero} value. It does not decide whether an arbitrary polynomial
has a zero among language outputs.

\subsection{Potential positivity (F34)}\label{sec:f34}
\entry{F34}{A.~Miasnikov and V.~Shpilrain}
Part (a) asks whether some automorphism makes a supplied word positive
in the fixed basis. Part (b) asks whether enlarging the ambient free
rank creates new potentially positive words; its negative answer is prior~\cite{clark-goldstein}.

\begin{theorem}
Potential positivity is decidable uniformly in finite rank.
\end{theorem}
\begin{proof}
An injection $\phi:F_r\to F_s$ reflects potential positivity. Indeed,
orient every edge of the based Stallings graph of $\phi(F_r)$ by its
positive ambient label. Choose a spanning tree and give each non-tree
edge this orientation in the resulting free basis. A positive closed
path stays positive when the tree is collapsed, so $\phi(w)$ positive
implies $w$ positive in a basis of its domain. No finite-index assumption
is needed.

Consequently $w$ is potentially positive exactly when there is an
endomorphism with $\phi(w)$ positive and $\det B_\phi\ne0$.
Encode $w(X_1,\ldots,X_r)=V$ with $V\in A^*$ and construct its
full relation~\eqref{eq:edt0l}. The determinant is a polynomial in
the separate counts of the $X_j$. Lemma~\ref{lem:polynomial} decides
whether it is identically zero. A nonzero output gives an injection
and the graph argument supplies a witnessing automorphism.
The identity is positive under the site's no-negative-letter convention;
an alternative nonempty convention changes only that boundary case.
\end{proof}

Rank-two decisions and the stability answer are credited to earlier
work in the \record{problems/F34/part-a-audit.md}{scope audit}.
The complete equation-to-language construction is imported and was not
implemented. The graph and polynomial stages and the language interfaces
were implemented and checked; that is a narrower assertion.
\evidence{F34}{part-a-proof.md}

\subsection{Translation equivalence (F38)}\label{sec:f38}
\entry{F38}{I.~Kapovich and P.~Schupp}
Part (a) asks whether $\clen{\alpha(u)}=\clen{\alpha(v)}$ for
all automorphisms can be decided. Part (c) replaces equality by
two-sided bounds on their ratio, uniform in the ambient automorphism.
Part (b), a substitution symmetry for translation-equivalent pairs,
is already answered affirmatively by Lee~\cite{lee-substitution}.

\begin{theorem}\label{prop:translation}
Translation equivalence is decidable uniformly in finite rank.
\end{theorem}
\begin{proof}
Kapovich--Levitt--Schupp--Shpilrain identify equality under all
automorphisms with equality under all injective endomorphisms
\cite[Corollary 1.4]{klss}. Thus the question is equivalent to
\[
 \det B_\phi\bigl(\clen{\phi(u)}-\clen{\phi(v)}\bigr)=0
 \quad\hbox{for every endomorphism }\phi.
\]
Nonzero determinant guarantees injectivity; conversely every
automorphism has determinant $\pm1$. Encode
$P^{-1}u(X)P=U$ and $Q^{-1}v(X)Q=V$, with $U,V$ cyclically
reduced, allowing the empty word. Every endomorphism extends to such
a tuple, and $|U|,|V|$ are exactly the two cyclic lengths. The displayed
expression is a polynomial of degree $r+1$ in separate letter counts.
Apply Lemma~\ref{lem:polynomial} to the full solution relation.
\end{proof}

For bounded equivalence write $H_w\le\Out(F_r)$ for the stabilizer
of the \emph{oriented} conjugacy class of $w$.

\begin{theorem}\label{prop:bounded}
For nonidentity $u,v$, the inequalities
$\clen{\alpha(v)}\le C\clen{\alpha(u)}$ for some $C$ and all
$\alpha\in\Aut(F_r)$ hold exactly when $H_u[v]$ is finite.
Hence bounded translation equivalence is equivalent to commensurability
of $H_u,H_v$, and is decidable.
\end{theorem}

Necessity follows because bounded cyclic length leaves finitely many
conjugacy classes. The proposed sufficiency uses a quantitative relative
shortening lemma, proved from Sela's \emph{graded} shortening theorem,
with a separate treatment of primitive powers in rank two. The exact
reduction and its normalization are in Appendix~\ref{app:shortening}.
An unqualified replacement of automorphisms by arbitrary injections
would be false for this bounded question.

The finite decision step is more elementary once the criterion is
granted. Whitehead peak reduction computes generators for $H_u,H_v$.
Let $I=\ker(\Out(F_r)\to\GL(r,\F_3))$. Handel--Mosher's
aperiodicity theorem says a periodic conjugacy class under an element
of $I$ is fixed~\cite[Theorem 4.1]{handel-mosher}. For a finitely
generated $H$, the orbit $H[w]$ is finite exactly when $H\cap I$
fixes $[w]$. Explore the finite mod-three image, retaining actual
automorphism representatives and the image of $[w]$ at each matrix.
Different classes at the same matrix certify an element of $H\cap I$
moving $[w]$ and therefore an infinite orbit. Otherwise the finite
graph closes consistently. Apply this to both orbits.

The code implements that finite procedure; it does not calculate a
shortening constant or validate the graded theorem. Rank-two decisions
are prior~\cite{lee,lee-translation}; the proposed additional range is higher rank.
For identity inputs, the ratio is undefined; its natural two-sided
inequality extension is true for two identities and false for exactly
one. Rank one is immediate.
\evidence{F38}{bounded-proof.md}
